%-------------------------------------------------------------------------------
% SECTION TITLE
%-------------------------------------------------------------------------------
\cvsection{Research Experience}

%-------------------------------------------------------------------------------
% CONTENT
%-------------------------------------------------------------------------------
\begin{cventries}

%---------------------------------------------------------
  \cventry
    {Research Assistant | Advisor: \href{https://pitt-photonics.github.io/}{Nathan Youngblood}}
    {University of Pittsburgh}
    {Pittsburgh, PA}
    {Sep 2019 - present}
    {
      \begin{cvitems}
        \item {Led experimental development of remote optical injection locking between a long-haul master laser and a local slave laser, enabling stable phase-coherent operation for interference-based signal processing, coherent detection, and optical-domain computing.}
        \item {Designed integrated photonic waveguides and components in KLayout, then validated fabricated chips through e-beam lithography workflows and optical characterization of coupling loss, interference contrast, and modulator response.}
        \item {Developed device- and system-level models in Ansys Lumerical, linking modulator efficiency, insertion loss, and coupling loss to locking range, SNR, and bandwidth; compared model predictions with measured system behavior.}
        \item {Built object-oriented Python automation with NI MAX and PyVISA/SCPI to coordinate lasers, SVAs, OSAs, and related instruments, creating reusable instrument and experiment classes for parameter sweeps, batch acquisition, and standardized analysis.}
        \item {Fabricated tapered fibers with embedded Fiber Bragg Gratings and designed fiber packaging for stress and temperature sensing; used COMSOL multiphysics to optimize taper geometry and thermo-mechanical response.}
        \item {Built a BOTDR setup for distributed temperature sensing and developed FBG/BOTDR calibration procedures, including temperature-wavelength mapping, sensitivity extraction, system-drift compensation, and regression-based analysis of uncertainty, noise, repeatability, and linearity.}
      \end{cvitems}
    }

%---------------------------------------------------------
  \cventry
    {Research Assistant | Advisor: \href{https://labsites.rochester.edu/zhang/}{Xi-Cheng Zhang}}
    {University of Rochester}
    {Rochester, NY}
    {Sep 2017 - Aug 2019}
    {
      \begin{cvitems}
        \item {Generated THz signals from liquid plasmas under varying temperature, polarity, and viscosity, quantifying how liquid properties influenced emission strength and spectral content.}
        \item {Used high-intensity ultrafast lasers to machine micron-scale apertures in syringe needles, controlling gas flow to produce acoustic tone spectra for gas-acoustic interaction studies.}
        \item {Implemented a digital cloaking prototype using commercially available lenslet arrays and MATLAB-based pixel remapping.}
        \item {Processed experimental data in time and frequency domains, using SNR, repeatability, sensitivity analysis, and regression fits to connect experimental parameters with sensing and THz performance.}
      \end{cvitems}
    }

\end{cventries}
